\section{Sources and Research Apparatus}

\begingroup\footnotesize

The entries below identify sources actually used. Web links were checked 12 July 2026. Biblical citations follow standard chapter-and-verse numbering. Latin liturgical incipits are checked against the 1962 typical-edition facsimile; English explanations in this study are original paraphrases unless a translation is expressly named.

\subsection{Controlling liturgical and magisterial sources}

\begin{itemize}
  \item \emph{Missale Romanum ex decreto SS. Concilii Tridentini restitutum Summorum Pontificum cura recognitum}, editio typica (1962), especially \emph{Ritus servandus in celebratione Missae}, \emph{Ordo Missae}, and \emph{Canon Missae}. \href{https://media.churchmusicassociation.org/pdf/missale62.pdf}{Public typical-edition facsimile}.
  \item John XXIII, apostolic letter \emph{Rubricarum instructum} (25 July 1960), nos. 1--5, promulgating the new Code of Rubrics. \href{https://www.vatican.va/content/john-xxiii/la/motu_proprio/documents/hf_j-xxiii_motu-proprio_19600725_rubricarum-instructum.html}{Holy See Latin text}.
  \item Council of Trent, Session XIII, decree on the Eucharist, especially chs. 1--5 and canons 1--4; Session XXII, doctrine and canons on the Sacrifice of the Mass, especially chs. 1--9 and canons 1--9.
  \item Pius V, apostolic constitution \emph{Quo primum} (14 July 1570), read with the title and contents of the promulgated Missal rather than as a claim that the Roman rite began in 1570.
  \item Pius XII, encyclical \emph{Mediator Dei} (20 November 1947), especially nos. 20--23, 47--59, 68--78, and 80--104. \href{https://www.vatican.va/content/pius-xii/en/encyclicals/documents/hf_p-xii_enc_20111947_mediator-dei.html}{Holy See English text}.
  \item Second Vatican Council, constitution \emph{Sacrosanctum Concilium} (4 December 1963), especially nos. 7, 14, and 47--58. \href{https://www.vatican.va/archive/hist_councils/ii_vatican_council/documents/vat-ii_const_19631204_sacrosanctum-concilium_en.html}{Holy See English text}.
  \item Second Vatican Council, dogmatic constitution \emph{Lumen gentium}, nos. 10--11; decree \emph{Presbyterorum ordinis}, nos. 2 and 5.
  \item \emph{Catechism of the Catholic Church}, nos. 1066--1209, 1322--1419, and 1544--1553. \href{https://www.vatican.va/archive/ENG0015/_INDEX.HTM}{Holy See English edition}.
\end{itemize}

\subsection{Scriptural, patristic, and early liturgical witnesses}

\begin{itemize}
  \item Genesis 14:18--20; Exodus 12; Isaiah 6:1--7; Malachi 1:11; Matthew 5:23--24; 6:9--13; 21:9; 26:26--29; Luke 22:14--20; 24:13--35; John 1:1--14; 6:22--71; Acts 2:42--47; 20:7; Romans 12:1--2; 1 Corinthians 10:16--21; 11:23--29; Hebrews 7--10; Revelation 4--5 and 19:6--9.
  \item \emph{Didache}, chs. 9--10 and 14.
  \item Justin Martyr, \emph{First Apology}, chs. 65--67.
  \item Irenaeus of Lyons, \emph{Against Heresies}, IV.17--18 and V.2.
  \item Tertullian, \emph{De corona} 3; Origen, \emph{Commentary on John} I.6; \emph{Apostolic Constitutions} VII.47.
  \item Cyprian of Carthage, \emph{Epistle} 63, on the mixed chalice, Christ's institution, and the unity of the people with their Lord; \emph{De dominica oratione} 8--9, 18, 31.
  \item The redactionally complex church order conventionally called the \emph{Apostolic Tradition}, used with the modern caution that date, authorship, and recoverable strata are disputed.
  \item Cyril of Jerusalem, \emph{Mystagogical Catechesis} V.
  \item Ambrose of Milan, \emph{De mysteriis}, chs. 8--9; \emph{De sacramentis}, IV.4--6 and V.
  \item Augustine of Hippo, Sermons 57, 227, 272, and 293; \emph{Enarrationes in Psalmos} 30, 42, and 85; \emph{Tractates on John} 1--2; \emph{City of God} X.6 and X.20.
  \item Gregory the Great, \emph{Register of Letters}, IX.26 (IX.12 in the NPNF arrangement), used as a bounded witness to Roman Kyrie, Alleluia, and Pater placement rather than proof that he fixed every later word.
  \item \emph{Liber pontificalis}, life of Sergius I, no.~86, for the Agnus Dei notice.
  \item \emph{Ordo Romanus I}, in Michel Andrieu, \emph{Les Ordines Romani du haut moyen âge}, vol.~II, \emph{Les textes (Ordines I--XIII)} (Louvain, 1948), as a witness to early medieval papal stational worship and differentiated ministry.
\end{itemize}

\subsection{Scholastic, medieval, and historical witnesses}

\begin{itemize}
  \item Amalarius of Metz, \emph{Liber officialis}; Innocent III, \emph{De sacro altaris mysterio}; William Durandus, \emph{Rationale divinorum officiorum}, book IV. These are witnesses to received medieval interpretation and must not automatically be treated as evidence of a ceremony's first cause.
  \item Thomas Aquinas, \emph{Summa theologiae} III, qq. 73--83, especially q. 75, aa. 1--4; q. 79; q. 82; and q. 83, aa. 1, 4--6.
  \item \emph{Missale Romanum} (Milan, 1474), and the typical editions promulgated in 1570, 1604, and 1634, used as major witnesses in the printed Missal's codification and controlled revision.
  \item Josef A. Jungmann, \emph{The Mass of the Roman Rite: Its Origins and Development}, trans. Francis A. Brunner, 2 vols. (New York: Benziger, 1951--1955), used as a major but not infallible historical synthesis.
  \item Eric Palazzo, \emph{A History of Liturgical Books from the Beginning to the Thirteenth Century}, trans. Madeleine Beaumont (Collegeville: Liturgical Press, 1998), for book types and the development of the plenary Missal.
  \item John F. Baldovin, \emph{The Urban Character of Christian Worship: The Origins, Development, and Meaning of Stational Liturgy} (Rome: Pontifical Oriental Institute, 1987), for the ecclesial and urban setting of early Roman stational worship.
  \item Uwe Michael Lang, \emph{The Roman Mass: From Early Christian Origins to Tridentine Reform} (Cambridge: Cambridge University Press, 2022), used as a recent historical control alongside primary witnesses.
\end{itemize}

\subsection{Thirty-three-step provenance}

\begin{itemize}
  \item John T. McMahon, \emph{Pray the Mass}, Part I, 13th revised ed. (Sydney: Pellegrini, 1959), with imprimatur of Archbishop Raymond Prendiville (June 1959). The booklet explicitly arranges thirty-three steps in association with Christ's earthly years. \href{https://www.ccwatershed.org/pdfs/pray-mass-booklet/download/}{Public scan}; \href{https://www.ccwatershed.org/2014/09/13/pray-mass-booklet-and-cards/}{bibliographic description}. The present study acknowledges this precedent but substitutes the comprehensive sequence documented in its repository inventory.
\end{itemize}

\begin{massbox}{Research status}
The 1962 Missal controls the received sequence. Patristic and medieval sources illuminate doctrinal and spiritual reception without being forced into a modern unit scheme. Historical dating is stated at the level supported by the evidence. Independent review by a liturgical historian, a patristics specialist, and a sacramental theologian remains outstanding.
\end{massbox}
\endgroup
